\section{Detection by Hecke modifications}\label{det:section}

We use the geometry of Section~\ref{geom:section} to prove two detection
statements. Over a field, the moving Hecke eigenrelations exclude every
nonnilpotent direction from singular support. Over a trait, they prevent
nearby cycles from vanishing when the generic support approaches the special
fiber. The second assertion will preserve nonzero objects through both
of the degenerations in the construction.

The argument adapts \cite[Theorems~4.1--4.2 and Lemmas~4.3--4.4]{CT}
to several ordinary or enhanced flags, their products, and their simultaneous
torus quotient. We reproduce the common local argument: one transverse
Hecke modification isolates a stalk in a proper pushforward, while a second
use of the modification recovers the original cut from an exact split
quadratic equation. The moving-leg microlocal method is related to
\cite[Sections~20.5--20.8 and Appendix~H]{AGKRRV}; here the precise two-fiber
input is Proposition~\ref{geom:transverse}.

Write $\cB$ for one of the stacks of Section~\ref{found:level-stacks},
and $\Lambda\subset T^*\cB$ for the generically nilpotent cone
of Section~\ref{geom:section}. For a smooth scheme chart $b:S\to\cB$,
write $\Lambda_S=b^\circ\Lambda\subset T^*S$. Products and quotients
have Hecke modifications on both open curves. Throughout this section the
base fields and characteristics are those of Section~\ref{geom:section}.

\begin{theorem}[Field detection]\label{det:field}
Let $F\in\Dlc(\cB)$ be a Hecke eigencomplex with lisse eigenvalues.
In the product and quotient cases assume the eigenrelations on both open
curves. Then, for every smooth finite-type chart $b:S\to\cB$,
\[
                           \SS(b^*F)\subset\Lambda_S.
\]
For this conclusion it is enough to have the individual one-leg
isomorphisms $\Hecke_V(F)\simeq F\boxtimes\mathcal V_V$ for all Satake
representations $V$, with each $\mathcal V_V$ lisse; their coherence is
not needed in this section.
\end{theorem}

For the second statement let $R$ be an excellent complete strictly
henselian discrete valuation ring with algebraically closed residue field
and $\ell$ invertible. Put $S_0=\Spec R$, with special point $s$, generic
point $\eta$, and a fixed geometric generic point $\bar\eta$.
Let $\cB\to S_0$ be smooth and locally of finite type, with special
fiber one of the preceding stacks. Assume a smooth relative curve of
legs $L\to S_0$ whose special fiber contains a nonempty open in each
curve on which modifications are allowed. The bounded Hecke maps over
input-and-leg and output-and-leg are representable and proper; the
exact-position maps are smooth, have the geometry of
Proposition~\ref{geom:transverse} on the special fiber, and carry the
standard open-stratum restriction of the Satake kernel. We use the
scheme or algebraic-space correspondence charts and split Schubert local
models described in Section~\ref{found:section}.

\begin{theorem}[Specialization detection]\label{det:specialization}
In this relative setup, let $F\in\Dlc(\cB_{\bar\eta})$ be a Hecke
eigencomplex with lisse eigenvalues on $L_{\bar\eta}$ satisfying the
lattice specialization condition of Proposition~\ref{found:nearby}.
If $\Psi F=0$, then on every smooth finite-type chart $D\to\cB$
the closure of the generic nonzero-stalk locus of $F|_{D_{\bar\eta}}$
does not meet $D_s$.
\end{theorem}

The closure in this statement can be computed after descending its finitely
many geometric generic components to a finite extension of $R$.  It does
not require a field of definition for $F$ itself.  All extensions below are
taken inside the fixed geometric generic field.  Nearby cycles always mean
geometric nearby cycles, without inertia invariants. For a relative chart
$D\to\cB$, the notation $\Lambda_D$ below means the smooth cotangent
pullback of the special-fiber cone to $T^*D_s$.

\subsection{A hypersurface cut and two Hecke modifications}

The common step in the two theorems is a calculation for one hypersurface.
We keep the two coefficient complexes distinct.  In the \emph{field case},
let $F$ satisfy the hypotheses of Theorem~\ref{det:field},
take a smooth chart $b:D\to\cB$, a closed point $d$, and a regular function
$f$ near $d$ with $f(d)=0$, and put $K=b^*F$.
Here $L$ is the open curve of allowed legs, or the disjoint union of
the two open curves in the product and quotient cases.  In the \emph{trait case}, let $F$
satisfy the hypotheses of Theorem~\ref{det:specialization}
and assume $\Psi F=0$.  Take a smooth chart $b:D\to\cB$ over $S_0$,
a closed point $d\in D_s$, and a regular function $f$ with $f(d)=0$;
put $K=b_{\bar\eta}^*F$.  Thus $K$ is defined on $D$ in the field case
and only on $D_{\bar\eta}$ in the trait case.

\begin{lemma}[Hypersurface calculation]\label{det:cut}
If the differential $df_d$ in the field case, respectively its relative
special-fiber differential in the trait case, lies outside $\Lambda_D$,
then
\[
      (\Phi_f K)_d=0
      \quad\text{in the field case},\qquad
      \bigl(\Psi(K|_{\{f=0\}_{\bar\eta}})\bigr)_d=0
      \quad\text{in the trait case}.
\]
Here $\Phi_f$ is the unshifted cone from restriction to geometric nearby
cycles for $f$.
\end{lemma}

\begin{proof}
The Hecke eigenrelation first gives a vanishing after proper pushforward.
The two nondegenerate zeros of Proposition~\ref{geom:transverse}
then have different roles: one isolates a single stalk in that pushforward,
and the other gives local coordinates in which the equation is the original
hypersurface equation plus a split quadratic form.  A projective compactification
will recover the cycles of the original hypersurface from those of the
stabilized one.

\smallskip\noindent
\emph{1. The modified hypersurface and the eigenrelation.}
If $df_d$ is nonzero on the tangent space relative to $b$, the map
$(b,f):D\to\cB\times\mathbb A^1$ is smooth near $d$ in the field case,
and $\{f=0\}\to\cB$ is smooth near $d$ in the trait case.  Smooth base
change proves the assertion.  We may therefore assume
$df_d=b^*A$ for a covector $A\notin\Lambda$ at $b(d)$.

Choose the modification $h$ supplied by
Proposition~\ref{geom:transverse}, at a leg $y$.  Write $H$ for its
exact-position Hecke stratum, with leg map $l:H\to L$, and write
\[
 p:H\longrightarrow\cB,\qquad
 o:H\longrightarrow\cB,\qquad
 q_H=(o,l):H\longrightarrow\cB\times L,
 \qquad \widetilde p=(p,l).
\]
Both $q_H$ and $\widetilde p$ are smooth of relative dimension $m$.  At $h$
the proposition gives
\begin{equation}\label{det:eq-leg}
                 p^*A=q_H^*(A',\xi),\qquad \xi\ne0.
\end{equation}
Use the two fibers named in that proposition:
\[
 H_{\mathrm{in}}=\widetilde p^{-1}(b(d),y),\qquad
 H_{\mathrm{out}}=q_H^{-1}(o(h),y).
\]
On $H_{\mathrm{in}}$, the restriction of $p^*A$ to the $q_H$-relative
tangent has a nondegenerate zero at $h$; on $H_{\mathrm{out}}$, the
restriction of $o^*A'$ to the $\widetilde p$-relative tangent does too.
Each fiber retains the identification with its fixed-side bundle.

In the trait case these are special-fiber statements; take the
corresponding relative Hecke stratum over $S_0$.  In the field case introduce
the completed local trait $S_0$ at $0$ in the function line, with
parameter $z$, extend all spaces
constantly, and impose $f=z$.  We will apply vanishing cycles over this
trait.  Thus the two cases have a common notation: $e=0$ in the trait
case and $e=z$ in the field case, and the operation to be computed is,
respectively, $\Psi$ and $\Phi$.

Take two copies $H_1,H_2$ of $H$ and form
\begin{equation}\label{det:eq-double}
 W_0=H_2\times_{p,\cB}D,
 \qquad
 P=H_1\times_{q_H,\cB\times L,q_H}W_0.
\end{equation}
Thus the two modifications have the same output bundle and leg, while
the input of the second lies in the chart $D$; Figure~\ref{det:figure}
records these two different roles.

\begin{figure}[htbp]
\centering
\begin{tikzpicture}[>=Stealth,baseline=(current bounding box.center)]
\node (D) at (0,0) {$D$};
\node (W) at (3,0) {$W_0$};
\node (P) at (3,1.8) {$P$};
\node (H) at (6.2,1.8) {$H_1$};
\node (O) at (6.2,0) {$\cB\times L$};
\node (B) at (9,1.8) {$\cB$};
\draw[->] (P) -- node[above] {$\operatorname{pr}_1$} (H);
\draw[->] (P) -- node[left] {$\operatorname{pr}_2$} (W);
\draw[->] (H) -- node[right] {$q_{H,1}$} (O);
\draw[->] (W) -- node[below] {$q_{H,2}$} (O);
\draw[->] (W) -- node[below] {\small second input} (D);
\draw[->] (H) -- node[above] {$p_1$} (B);
\end{tikzpicture}
\caption{The cartesian square of two modifications with a common output
and leg. The cut is imposed through the left-hand map to $D$; the
coefficient is pulled from the first input at the upper right.
The diagonal section $W_0\to P$ makes the two modifications equal.}
\label{det:figure}
\end{figure}

Let $c:P\to D$ be the projection through $W_0$.  Set
$W=\{f=e\}\subset W_0$ and $P_c=\{f(c)=e\}\subset P$.
The distinguished points are $w=(h,d)$ in $W_s$ and
$a=(h,h,d)$ in $(P_c)_s$.  Products here are over the trait when one is
present.

The output-bundle map $W\to\cB$ is smooth near $w$.  Indeed,
$W_0\to\cB\times L$ is smooth, and
\eqref{det:eq-leg} says that the cut differential, on directions
relative to the output bundle, has the nonzero leg component $\xi$.
The eigenisomorphism therefore implies that the cycles of the pulled-back
Hecke transform vanish at $w$.  In the trait case this is smooth base
change from $\Psi F=0$, followed by the lisse tensor compatibility.  In
the field case $W\to\cB_{S_0}$ is smooth over the function trait, where
$\cB_{S_0}$ is the constant stack.  Pullback of the constant family with
coefficient $F$ therefore has zero vanishing cycles, and tensoring by
the lisse eigenvalue preserves this assertion.  Only the smooth map to the bundle stack enters this argument.

\smallskip\noindent
\emph{2. Isolating one stalk in the proper pushforward.}
Replace $H_1$ by its Schubert bound $\overline H_1$ and let
$\IC_\lambda$ be the simple Satake kernel of this bound. Set
\[
 \overline P_c=\overline H_1\times_{\cB\times L}W,
 \qquad \pi_c:\overline P_c\longrightarrow W.
\]
This map is proper. Let $\mathcal J$ be the pullback to $\overline P_c$
of the Hecke integrand $F\,\widetilde\boxtimes\,\IC_\lambda$ on
$\overline H_1$, with the twisted-product convention of
\eqref{found:hecke}. In the field case extend the integrand constantly
in $z$ before pulling back to the cut; thus $\mathcal J$ is defined on
the whole $z$-family. In the trait case it is defined on the geometric
generic fiber. Proper base change identifies $R\pi_{c,*}\mathcal J$
with the pullback to $W$ of the Hecke transform used in Step~1.

Let $T=\pi_c^{-1}(w)$ and let $C$ be the restriction to $T$ of
$\Phi\mathcal J$ in the field case, or of $\Psi\mathcal J$ in the
trait case. Proper compatibility of cycles and the vanishing in Step~1
give
\begin{equation}\label{det:eq-proper-zero}
                            R\Gamma(T,C)=0.
\end{equation}
On the open stratum $P_c\subset\overline P_c$ the kernel is
$E[m](m/2)$. We will isolate the contribution of $a$ there.

The open part of $T$ is exactly $H_{\mathrm{out}}$ for $H_1$,
since $w$ fixes $H_2=h$ and $c=d$. Thus \eqref{det:eq-leg} at the
fixed second modification identifies $d(f\circ c)$ along this whole
fiber with the pullback of $(A',\xi)$ along $q_{H,1}$.
Before cutting, $P\to\cB$ by the first input bundle is smooth. At a
point of $T$ in its open-orbit neighborhood, failure of smoothness after
cutting would mean that $d(f\circ c)$ comes from the cotangent space of
that first input bundle. Since $P\to H_1$ is smooth, its
cotangent map is injective; thus failure would already give on $H_1$
an equality
\[
                         q_{H,1}^*(A',\xi)=p_1^*A_1
\]
for some $A_1$.  On $\widetilde p_1$-relative tangents this forces
$o_1^*A'$ to vanish.  The nondegenerate zero on $H_{\mathrm{out}}$
shows that $a$ is isolated among such points of $T$.  Consequently $C$
vanishes on a punctured neighborhood of $a$ in $T$: there the cut map
to the first input stack is smooth and the open-orbit Satake kernel is
an invertible constant shift and twist.
This smoothness is over the trait.  In the field case the integrand
therefore pulls back the constant $F$-family on $\cB_{S_0}$, so it has zero
$\Phi$, as required; the trait case uses $\Psi F=0$.

Here is the precise consequence of this isolation.  If a constructible
complex on a proper scheme vanishes on a punctured neighborhood of a
closed point, its contribution at that point is a direct summand of its
global cohomology.  To see this, let $j:T\setminus\{a\}\to T$ and
$i:\{a\}\to T$.  The restriction $j^*C$ is zero near $a$, so
$(Rj_*j^*C)_a=0$.  The localization triangle consequently splits as the
direct sum of $i_*C_a$ and $Rj_*j^*C$; both connecting morphisms vanish
by these disjoint supports.  The same localization argument applies
to algebraic spaces.  Equation
\eqref{det:eq-proper-zero} now gives
\begin{equation}\label{det:eq-vertex-zero}
                                  C_a=0.
\end{equation}

\smallskip\noindent
\emph{3. An exact split quadratic equation at the isolated point.}
We identify the local equation at this isolated point exactly.  The
diagonal $W_0\to P$, given by making the two modifications equal, is a
section of the smooth map $P\to W_0$.  \'Etale locally at $a$, there is a
map $r:P\to D$ lifting the first-input map to $\cB$, whose restriction
to this diagonal is the original chart point, including its
identification with the first input.  Indeed, the smooth morphism
$P\times_{\cB}D\to P$ has that section on the diagonal.  To extend it,
choose relative smooth coordinates, giving an \'etale map from a
neighborhood of its value to $\mathbb A^\delta_P$, where $\delta$ is
the relative dimension of $D\to\cB$ at $d$.  The coordinates of the
prescribed section are functions on the diagonal; lift these functions
to a neighborhood in $P$ and take their graph in $\mathbb A^\delta_P$.
Pullback of the \'etale map along this graph gives an \'etale neighborhood
of $P$ with a lift to $P\times_{\cB}D$.  Retain the branch carrying
the prescribed section along the diagonal.  The resulting projection
to $D$ is $r$, and the fiber product retains the specified stack
isomorphism between $b\circ r$ and the first input.
Thus $c$ records the second input everywhere, whereas $r$ records the
first input in this \'etale neighborhood; the two agree on the diagonal.
The diagonal, being a section of the smooth morphism $P\to W_0$ of
relative dimension $m$, is a regular immersion of codimension $m$.
Choose relative coordinates $v_1,\ldots,v_m$ generating its ideal.
Since $r=c$ on it, there are functions $u_i$ with
\begin{equation}\label{det:eq-exact-uv}
                         f(c)-f(r)=\sum_{i=1}^m u_i v_i.
\end{equation}
In particular, the subsequent change of coordinates will give the exact
cut equation.

On the diagonal fiber with $c=d$ and leg $y$, which is the fixed-input,
fixed-leg fiber of $H$ through $h$, the conormal coefficients
$u_i$ are the negatives of the restriction of $p^*A$ to the
$q_H$-relative tangent, in the frame dual to the $v_i$.  In fact variation
of the first modification with the second fixed has $dc=0$, whereas
$b\circ r$ is its first input, so differentiating
\eqref{det:eq-exact-uv} gives precisely this restriction.  This
calculation is independent of the chart-vertical part of $dr$:
$df_d=b^*A$ annihilates $\ker(db_d)$.
The description holds along that entire diagonal fiber.  The
nondegenerate zero on $H_{\mathrm{in}}$ shows that $u_i(a)=0$ and that the
derivatives of the $u_i$ on that fiber form a basis.  The map
$W_0\to D\times L$ is the smooth base change of $\widetilde p$,
of relative dimension $m$.  On the diagonal,
$r=c$ and the leg provide the other coordinates; the $v_i$ provide its
normal coordinates.  Hence
\begin{equation}\label{det:eq-uv-chart}
 (r,l,u_1,\ldots,u_m,v_1,\ldots,v_m):
 P\longrightarrow D\times L\times\mathbb A^{2m}
\end{equation}
is \'etale near $a$, with the evident relative interpretation over $S_0$.
On this neighborhood of the cut,
$\mathcal J|_{P_c}\simeq r^*K[m](m/2)$. Thus
\eqref{det:eq-vertex-zero} says that the cycles of the pulled-back $K$
vanish at the vertex of
\begin{equation}\label{det:eq-affine-quadric}
                         f+\sum_i u_i v_i=e.
\end{equation}
For $m=0$ this is already the required assertion, after removing the
smooth leg factor.

\smallskip\noindent
\emph{4. Recovering the original hypersurface.}
We have proved vanishing after adjoining the split quadratic variables.
To remove them, we use a proper family whose extra cohomology is
supported precisely on the original cut $\{f=e\}$.
For $m>0$, put $B=D\times L$ and compactify
\eqref{det:eq-affine-quadric} to the proper quadric family
\begin{equation}\label{det:eq-projective-quadric}
 \pi:Q=\left\{\sum_{i=1}^m u_i v_i+(f-e)w^2=0\right\}
               \subset\mathbb P^{2m}_B\longrightarrow B.
\end{equation}
In the field case $B$ also carries the constant extension in $z$.
Every point above $(d,y)$ other than the vertex $[0:\cdots:0:1]$ is
smooth for $\pi$, since some $u_i$ or $v_i$ is nonzero.  The cycles of
$\pi^*K$ there vanish by smooth base change; they vanish at the vertex
by \eqref{det:eq-uv-chart}.  Proper compatibility therefore gives
zero cycles for $R\pi_*\pi^*K$ at $(d,y)$.

The split-quadric calculation in
\cite[proof of Lemma~4.3]{CT} gives the precise extra term we need.
Write $i:Z=\{f=e\}\hookrightarrow B$. Hyperplane powers give a triangle
\begin{equation}\label{det:eq-quadric-triangle}
 \bigoplus_{a=0}^{2m-1}E_B(-a)[-2a]\longrightarrow R\pi_*E_Q
 \longrightarrow i_*E_Z(-m)[-2m]\xrightarrow{+1}.
\end{equation}
Here the first arrow is an isomorphism outside $Z$: a smooth split
quadric of dimension $2m-1$ has the same $E$-cohomology as
$\mathbb P^{2m-1}$. Over $Z$ the family is the constant projective cone
on the split quadric of dimension $2m-2$. Its cohomology has one
additional copy of $E(-m)$ in degree $2m$, beyond the hyperplane
powers. Proper base change identifies the restriction of the cone of
the first arrow with this constant complex on $Z$; localization
therefore identifies the cone as the last term displayed. This explains
why the triangle is a statement about sheaves, not merely fiber
cohomology dimensions.

Tensor \eqref{det:eq-quadric-triangle} with $K$ and use the
projection formula.  In the trait case take this triangle on the
geometric generic fiber.  The cycles of its first term vanish at
$(d,y)$ because $\Psi K=0$; those of its middle term vanish by the
proper calculation above.  Its last term therefore has zero nearby
cycles there, and proper compatibility for $i$ proves the claimed
vanishing on $\{f=0\}$.  In the field case use the triangle on the
whole $z$-family.  Its first term has zero vanishing cycles because
$K$ is pulled from the constant family $D$.  The identical argument
with $\Phi$ proves $(\Phi_f K)_d=0$.  Smooth base change removes the
leg factor in both cases.  The sheaf operations used throughout are
those in Proposition~\ref{found:nearby}.
\end{proof}

\subsection{Nilpotent singular support}

\begin{proof}[Proof of Theorem~\ref{det:field}]
Apply Lemma~\ref{det:cut} on every smooth chart and
every open subchart.  It says that every function whose differential
at a point is outside $\Lambda_D$ is locally acyclic there relative to
$b^*F$.  The cone includes the zero section, so these are precisely the
nonzero directions that must be excluded.  Weak singular support is
the closure of the differentials of functions with nonzero vanishing
cycles; over an infinite field closed-point tests suffice by
constructibility.  Barrett proves this description for the present
$E$-coefficients and proves that weak singular support equals full
singular support \cite[\S\S1.4--1.5, Theorem~(vii)]{Barrett}, extending
Beilinson's torsion-coefficient theorem \cite[\S1.5]{Beilinson}.
The hypersurface calculation therefore gives
$\SS(b^*F)\subset\Lambda_D$.  The same calculation on charts with
additional smooth parameters is available, so the conclusion is
compatible with the smooth-chart definition of singular support on
the stack.
\end{proof}

\subsection{Simultaneous cuts and nonvanishing of specialization}

To prove specialization detection, we must pass from hypersurfaces to
enough simultaneous cuts to make the support finite over the trait.
Successive hypersurface applications would require information about the
singular support of intermediate restrictions.  The following incidence
argument makes all the cuts at once.

\begin{lemma}[Simultaneous cuts]\label{det:simultaneous}
In the trait case, assume $\Psi F=0$.  Let $f_1,\ldots,f_r$ vanish at
$d\in D_s$ and suppose their special-fiber differentials span an
$r$-dimensional subspace $V\subset T_d^*D_s$ with
$V\cap(\Lambda_D)_d=\{0\}$.  Then nearby cycles of
$K|_{\{f_1=\cdots=f_r=0\}_{\bar\eta}}$ vanish at $d$.
The statement for $r=0$ is $\Psi K=0$.
\end{lemma}

\begin{proof}
The case $r=1$ is Lemma~\ref{det:cut}.  For $r>1$
consider the proper incidence projection
\[
 \pi:J=\{(t,[a_1:\cdots:a_r])\in D\times\mathbb P^{r-1}:
                   \textstyle\sum_i a_i f_i(t)=0\}\longrightarrow D.
\]
At $(d,[a])$ the differential of the defining function on a standard
projective chart is $\sum_i a_i df_i$; derivatives in the projective
directions are zero since all $f_i(d)$ vanish.  It is nonzero and lies
outside the pulled-back nilpotent cone.  The map
$D\times\mathbb A^{r-1}\to\cB$ on each such chart is smooth, so
Lemma~\ref{det:cut} proves that the nearby cycles of
$\pi^*K$ vanish on the whole fiber $\mathbb P^{r-1}$ over $d$.
Proper compatibility gives $(\Psi R\pi_*\pi^*K)_d=0$.

Let $i:Z=\{f_1=\cdots=f_r=0\}\hookrightarrow D$.  On the geometric
generic fiber hyperplane powers give the triangle
\[
 \bigoplus_{a=0}^{r-2}E_D(-a)[-2a]\longrightarrow R\pi_*E_J
 \longrightarrow i_*E_Z(-(r-1))[-2(r-1)]\xrightarrow{+1}.
\]
Indeed, over $D\setminus Z$ the incidence is a projective bundle of
relative dimension $r-2$, so the first arrow is an isomorphism there.
Over $Z$ it is the product $Z\times\mathbb P^{r-1}$, and the cone has
only its constant top cohomology sheaf.  Proper base change and
localization identify the cone as written, not merely its stalk
dimensions.  Tensor by $K$.  The first term has zero nearby cycles by
$\Psi K=0$, and the middle term by the preceding proper calculation.
Proper compatibility for $i$ proves the lemma.
\end{proof}

\begin{proof}[Proof of Theorem~\ref{det:specialization}]
Suppose that the closure of the nonzero-stalk locus meets $D_s$.
Shrink to an affine chart smooth of constant relative dimension $N$ over
$S_0$.  The
closure of that locus in $D_{\bar\eta}$ has finitely many irreducible
components $Z_{\bar\eta,\alpha}$.  Constructibility gives, on each
component, a dense open subset where the stalk of $K$ is nonzero.
After a finite extension of the trait, descend the components and the
proper closed complements of these opens.  We only descend these finite
algebraic data, and continue to use $K$ over $\bar\eta$.
Let $Z_\alpha\subset D$ and $E_\alpha\subset Z_\alpha$ be their
horizontal closures.  Components not meeting $D_s$ may be removed.

Write $j$ for the largest dimension of a generic component whose
closure meets $D_s$, and choose an irreducible component $Z_0$ of the
reduced special fiber of such a closure.  A horizontal integral
finite-type scheme over this excellent trait, with generic dimension
$j$, has every nonempty special-fiber component of dimension $j$;
its local dimension at a special-fiber closed point is $j+1$.
For these two assertions see, respectively,
\cite[Tags 0B2J and 02JU]{Stacks}.  The same dimension calculation
shows that the special fibers of the horizontal exceptional closures
$E_\alpha$ have dimension at most $j-1$ whenever the corresponding
generic component has dimension at most $j$.

Choose a general closed point $d$ of $Z_0$.  We may arrange that $Z_0$
is smooth at $d$, that $d$ is in none of the exceptional closures,
and that every reduced special-fiber component of the support closure
through $d$ equals $Z_0$ near $d$.  Distinct horizontal components may
have this same special component; this causes no difficulty.  By
Proposition~\ref{geom:dimension},
\[
                  \dim\Lambda_D\leq N.
\]
After shrinking a dense open in $Z_0$, the fiber dimension theorem
therefore gives
\begin{equation}\label{det:eq-cone-fiber}
                    \dim(\Lambda_D)_d\leq N-j.
\end{equation}

Choose a $j$-plane $V\subset T_d^*D_s$ which avoids
$(\Lambda_D)_d\setminus\{0\}$ and whose restriction to $T_dZ_0$
is an isomorphism onto $T_d^*Z_0$.  Both are nonempty open conditions
on the Grassmannian.  When $j>0$, for the first condition projectivize the cone in
\eqref{det:eq-cone-fiber}: it has dimension at most $N-j-1$,
and a general $\mathbb P^{j-1}$ in $\mathbb P^{N-1}$ is disjoint
from it by the elementary incidence dimension count.  The second is
the usual complement-to-a-fixed-subspace condition.  The base field
is infinite, so the two opens meet.  If $j=0$, take $V=0$ and use no
cutting functions.  Lift a basis of $V$ to
differentials of regular functions $f_1,\ldots,f_j$ vanishing at $d$.
Their restrictions are parameters on the smooth scheme $Z_0$.
Put $Y=\{f_1=\cdots=f_j=0\}\subset D$.
Lemma~\ref{det:simultaneous} says
\begin{equation}\label{det:eq-final-zero}
                         (\Psi(K|_{Y_{\bar\eta}}))_d=0.
\end{equation}

We finish by showing directly that this stalk is nonzero.  Near $d$,
the special fiber of the closed support closure intersected with $Y$
is zero-dimensional.  On any of its selected horizontal components
$Z_\alpha$ of generic dimension $j$ through $d$, the local ring has
dimension $j+1$, so its quotient by the $j$ functions has dimension
at least one \cite[Tag 0B52]{Stacks}.  Its further quotient by the uniformizer is
zero-dimensional.  Thus this intersection has a generization of $d$
in the generic fiber.  Such a generization lies outside $E_\alpha$,
since $d$ was chosen outside its closure, and $K$ has a nonzero
geometric stalk there.

For completeness, the passage from this generization to nearby cycles
is local and finite.  Let $Z\subset Y$ be the reduced intersection
$Y\cap\bigcup_\alpha Z_\alpha$.  Its generic fiber contains the
nonzero-stalk locus of $K|_{Y_{\bar\eta}}$, the preceding generization
lies in $Z$, and $d$ is isolated in $Z_s$.
The morphism $Z\to S_0$ is quasi-finite at $d$.
Take an affine neighborhood in $Y$ and use the henselian decomposition
of its closed support: the local ring of $Z$ at $d$ is a finite local
$R$-algebra and defines an open-and-closed piece $T\subset Z$
\cite[Tag 04GH(3) and its proof]{Stacks}.  Its unique special-fiber
point is $d$.  Remove the closed complement $Z\setminus T$ from the
ambient neighborhood.  On the resulting neighborhood, which has the
same nearby-cycle stalk at $d$, the support closure is the finite
$S_0$-scheme $T$.

The generic restriction of $K$ is the extension by zero from this
finite closed support.  Proper base change now identifies its
nearby-cycle stalk at $d$ with
\[
               \bigoplus_{z\in T_{\bar\eta}}
                             (K|_{Y_{\bar\eta}})_z.
\]
This is a finite direct sum of complexes of $E$-vector spaces, at
least one of which is nonzero.  Hence the sum is nonzero, contradicting
\eqref{det:eq-final-zero}.  The argument includes $j=0$, when
there are no cutting functions and the original support already has
the required finite local piece.  This proves the theorem.
\end{proof}
